\documentclass[aspectratio=169]{beamer}

% ---------- Russian + math ----------
\usepackage[T2A]{fontenc}
\usepackage[utf8]{inputenc}
\usepackage[russian]{babel}
\usepackage{amsmath, amssymb, mathtools}
\usepackage{booktabs}
\usepackage{array}
\usepackage{hyperref}
\usepackage{graphicx}
\usepackage{xcolor}
\usepackage{listings}

% ---------- Beamer theme ----------
\usetheme{Madrid}
\usecolortheme{default}
\setbeamertemplate{navigation symbols}{}

% ---------- Listings (Python) ----------
\lstdefinestyle{py}{
  language=Python,
  basicstyle=\ttfamily\small,
  keywordstyle=\bfseries\color{blue!60!black},
  commentstyle=\itshape\color{gray!70!black},
  showstringspaces=false,
  columns=fullflexible,
  frame=single,
  framerule=0.4pt,
  breaklines=true,
  tabsize=2
}
\lstset{style=py}

% ---------- Title ----------
\title{Введение в машинное обучение}
\subtitle{Лекция: Кластеризация и снижение размерности (SVD, PCA)}
\author{Лазар В.~И.\ и Козлова Е.~Р.}

% ---------- Helpers ----------
\newcommand{\term}[1]{\textbf{#1}}
\newcommand{\R}{\mathbb{R}}

\begin{document}

% ---------- Title slide ----------
\begin{frame}
	\titlepage
\end{frame}

% ---------- Plan ----------
\begin{frame}{План занятия}
	\begin{enumerate}
		\item Что такое «обучение без учителя».
		\item \term{Кластеризация}: идеи и алгоритмы
		      \begin{itemize}
			      \item K-Means
			      \item DBSCAN
			      \item Иерархическая (агломеративная) кластеризация
		      \end{itemize}
		\item \term{Снижение размерности}:
		      \begin{itemize}
			      \item SVD — сингулярное разложение
			      \item PCA — метод главных компонент
		      \end{itemize}
		\item Практические советы, мини-упражнения, ответы на вопросы.
	\end{enumerate}
\end{frame}

% ---------- Intro: Unsupervised ----------
\begin{frame}{Обучение без учителя (интуиция)}
	\begin{itemize}
		\item Даны только объекты (точки, картинки, тексты), \emph{без меток ответов}.
		\item Хотим: \term{найти структуру} в данных.
		\item Две базовые задачи:
		      \begin{enumerate}
			      \item \term{Кластеризация} — разбить объекты на «похожие группы».
			      \item \term{Снижение размерности} — «упростить» данные, сохранив самое важное.
		      \end{enumerate}
	\end{itemize}
\end{frame}

% =======================
% K-MEANS
% =======================
\begin{frame}{Кластеризация: K-Means — идея}
	\begin{itemize}
		\item Выбираем число групп \(K\).
		\item Каждая группа описывается центром (\term{центроидом}).
		\item Алгоритм стремится «поставить» центры так, чтобы точки были ближе к своим центрам, чем к чужим.
		\item Мера близости — обычно евклидово расстояние.
	\end{itemize}
\end{frame}

\begin{frame}{K-Means — пошаговый алгоритм}
	\begin{enumerate}
		\item \term{Инициализация}: выбрать \(K\) начальных центров (лучше \(\text{k-means++}\)).
		\item \term{Шаг присвоения}: каждую точку отнести к ближайшему центру.
		\item \term{Шаг пересчёта}: для каждого кластера посчитать новый центр как среднее своих точек.
		\item Повторять шаги 2–3, пока центры почти не меняются (\term{сходимость}).
	\end{enumerate}
	\vspace{0.5em}
	\[
		\text{Целевая функция: }
		\min_{\text{кластеры}} \sum_{i=1}^{n} \left\|x_i - \mu_{c(i)}\right\|^2
	\]
	где \(x_i\) — точка, \(\mu_{c(i)}\) — центр её кластера.
\end{frame}

\begin{frame}{K-Means — плюсы и минусы}
	\begin{columns}[T]
		\begin{column}{0.52\textwidth}
			\textbf{Плюсы}
			\begin{itemize}
				\item Простой и быстрый.
				\item Понятная геометрическая идея.
				\item Хорош для «круглых» кластеров схожего размера.
			\end{itemize}
		\end{column}
		\begin{column}{0.48\textwidth}
			\textbf{Минусы}
			\begin{itemize}
				\item Нужно знать \(K\) заранее.
				\item Плохо с «некруглыми» формами и разной плотностью.
				\item Чувствителен к масштабам признаков и выбросам.
			\end{itemize}
		\end{column}
	\end{columns}
	\vspace{0.5em}
	\textbf{Практика:} перед K-Means часто делают стандартизацию признаков.
\end{frame}

% =======================
% DBSCAN
% =======================
\begin{frame}{Кластеризация: DBSCAN — идея}
	\begin{itemize}
		\item \term{Density-Based Spatial Clustering of Applications with Noise}.
		\item Вместо «круглых» групп ищем \term{области высокой плотности}.
		\item Параметры:
		      \begin{itemize}
			      \item \(\varepsilon\) (eps) — «радиус соседства»;
			      \item \(\text{min\_samples}\) — сколько соседей нужно, чтобы точка считалась «ядром».
		      \end{itemize}
		\item Типы точек: \term{ядро}, \term{пограничная}, \term{шум}.
		\item Число кластеров получается автоматически.
	\end{itemize}
\end{frame}

\begin{frame}{DBSCAN — плюсы и минусы}
	\begin{columns}[T]
		\begin{column}{0.52\textwidth}
			\textbf{Плюсы}
			\begin{itemize}
				\item Находит \emph{кластеры произвольной формы}.
				\item Не требует знать число кластеров.
				\item Устойчив к выбросам (помечает как шум).
			\end{itemize}
		\end{column}
		\begin{column}{0.48\textwidth}
			\textbf{Минусы}
			\begin{itemize}
				\item Нужно подбирать \(\varepsilon\) и \(\text{min\_samples}\).
				\item Трудно при сильно неоднородной плотности.
				\item Чувствителен к масштабу признаков.
			\end{itemize}
		\end{column}
	\end{columns}
\end{frame}

% =======================
% Agglomerative
% =======================
\begin{frame}{Иерархическая (агломеративная) кластеризация — идея}
	\begin{itemize}
		\item Стартуем: каждый объект — свой кластер.
		\item На каждом шаге \term{сливаем} две наиболее близкие группы.
		\item Выбор «близости» групп — это \term{связь (linkage)}:
		      \begin{itemize}
			      \item \term{single} — ближайшие точки;
			      \item \term{complete} — самые дальние точки;
			      \item \term{average} — среднее расстояние;
			      \item \term{Ward} — минимизируем рост дисперсии.
		      \end{itemize}
		\item Получаем \term{дендрограмму}; «отрезаем» на нужной высоте — получаем кластеры.
	\end{itemize}
\end{frame}

\begin{frame}{Иерархическая кластеризация — плюсы и минусы}
	\begin{columns}[T]
		\begin{column}{0.52\textwidth}
			\textbf{Плюсы}
			\begin{itemize}
				\item Дает \emph{иерархию} — разные уровни детализации.
				\item Не обязательно заранее задавать \(K\).
				\item Гибкость выбора метрики и связей.
			\end{itemize}
		\end{column}
		\begin{column}{0.48\textwidth}
			\textbf{Минусы}
			\begin{itemize}
				\item Память/время \(\sim n^2\) — трудно для очень больших данных.
				\item Ранние «неудачные» слияния уже не откатить.
			\end{itemize}
		\end{column}
	\end{columns}
\end{frame}

% =======================
% Comparison table
% =======================
\begin{frame}{Коротко: что выбрать?}
	\small
	\begin{tabular}{@{}p{3.1cm} p{3.3cm} p{3.3cm} p{3.3cm}@{}}
		\toprule
		                  & \textbf{K-Means}             & \textbf{DBSCAN}               & \textbf{Иерархическая}             \\
		\midrule
		Форма кластеров   & «Круглые», схожие по размеру & Любые формы, плотностные      & Любые, зависит от метрики/связи    \\
		Заранее \(K\)?    & Да                           & Нет                           & Необязательно (режем дендрограмму) \\
		Выбросы           & Чувствителен                 & Устойчив (шум = -1)           & Зависит от метрики                 \\
		Гиперпараметры    & \(K\), инициализация         & \(\varepsilon\), min\_samples & Связь (linkage), «высота среза»    \\
		Масштаб признаков & Важно                        & Важно                         & Важно                              \\
		Сложность         & Быстрый                      & Средний                       & \(O(n^2)\) по памяти/времени       \\
		\bottomrule
	\end{tabular}
\end{frame}

% =======================
% Dimensionality reduction
% =======================
\begin{frame}{Снижение размерности — зачем?}
	\begin{itemize}
		\item Упростить: оставить «главное», убрать шум/избыточность.
		\item \term{Визуализация}: проецировать высокие измерения на 2D/3D.
		\item \term{Скорость}: меньше признаков — быстрее обучение.
		\item \term{Сжатие}: хранить данные компактнее.
	\end{itemize}
\end{frame}

% ---------- SVD ----------
\begin{frame}{SVD — сингулярное разложение (идея)}
	\begin{itemize}
		\item Для матрицы \(X \in \R^{n\times d}\):
		      \[
			      X = U\,\Sigma\,V^{\top},
		      \]
		      где \(U \in \R^{n\times n}\), \(V \in \R^{d\times d}\) — \term{ортогональные} матрицы,
		      \(\Sigma\) — диагональная с неотрицательными \term{сингулярными значениями} \(\sigma_1 \ge \sigma_2 \ge \dots \ge 0\).
		\item Геометрия: \(V\) — «направления» в признаковом пространстве, \(\Sigma\) — «силы» по этим направлениям, \(U\) — координаты точек в новых осях.
		\item Если оставить только первые \(k\) наибольших \(\sigma\), получаем \term{лучшую аппроксимацию} ранга \(k\) (в смысле наименьших квадратов).
	\end{itemize}
\end{frame}

% ---------- PCA ----------
\begin{frame}{PCA — метод главных компонент (интуиция)}
	\begin{itemize}
		\item Ищем новые оси, вдоль которых \term{разброс данных (дисперсия)} максимален.
		\item 1-я главная компонента — направление наибольшей дисперсии, 2-я — следующей и т.д.
		\item Проекция на первые \(k\) компонент сохраняет «больше всего информации» в смысле дисперсии.
	\end{itemize}
\end{frame}

\begin{frame}{PCA — пошаговый алгоритм}
	\begin{enumerate}
		\item \term{Центрирование}: вычесть среднее из каждого признака.
		\item (Часто) \term{Стандартизация}: привести признаки к одному масштабу.
		\item Посчитать ковариационную матрицу \(C = \frac{1}{n-1} X^\top X\) (для центрированного \(X\)).
		\item Найти собственные векторы/значения \(C\); отсортировать по убыванию.
		\item Взять первые \(k\) векторов, спроецировать: \(Z = X W_k\).
	\end{enumerate}
	\vspace{0.5em}
	\textbf{Связь с SVD:} если \(X = U\Sigma V^\top\) и \(X\) центрирована, то столбцы \(V\) — направления PCA, а \(\Sigma^2/(n-1)\) — дисперсии компонент.
\end{frame}

% ---------- SVD code ----------

% =======================
% Practical tips
% =======================
\begin{frame}{Практические советы}
	\begin{itemize}
		\item \textbf{Масштабируйте} признаки (StandardScaler), особенно перед K-Means/PCA.
		\item Для K-Means подберите \(K\): \term{локоть} (elbow) или \term{силуэт}.
		\item Для DBSCAN подберите \(\varepsilon\) по графику k-ближайших расстояний.
		\item В PCA смотрите на \term{долю объяснённой дисперсии}; часто берут 90–95\%.
		\item Удаляйте \term{явные выбросы} или используйте устойчивые алгоритмы (DBSCAN).
	\end{itemize}
\end{frame}

\end{document}
